\subsection{Contagem}
\begin{itemize}[leftmargin=*,label={}]
		%% ----------------------------------------------------- %%
		\ifextended
	\item \textbf{Coeficiente binomial.}
		\[
			\binom{n}{k}=\frac{n!}{k!(n-k)!}
		\]
		Número de formas de escolher $k$ elementos entre $n$ tipos, com repetição.
		\fi

		%% ----------------------------------------------------- %%
	\item \textbf{Estrelas e Barras.}
		O número de soluções inteiras não negativas de
		\[
			x_1+\cdots+x_k=n
		\]
		é
		\[
			\binom{n+k-1}{k-1}.
		\]
		O número de soluções inteiras positivas é
		\[
			\binom{n-1}{k-1}.
		\]

		%% ----------------------------------------------------- %%
		\ifextended
	\item \textbf{Estrelas e barras com cotas inferiores.}
		O número de soluções inteiras de
		\[
			x_1+\cdots+x_n = S,\qquad x_i\ge l_i
		\]
		é
		\[
			\binom{S-\sum l_i+n-1}{n-1},
		\]
		após a substituição $y_i=x_i-l_i$.
		\fi

		%% ----------------------------------------------------- %%
	\item \textbf{Número de Catalan.}
		\[
			C_n=\frac{1}{n+1}\binom{2n}{n}
			=\binom{2n}{n}-\binom{2n}{n+1}
		\]
		\ifextended
		conta sequências balanceadas de parênteses de tamanho $2n$, formas de árvores binárias enraizadas com $n$ nós internos, triangulações de um $(n+2)$-ágono, matchings perfeitos sem cruzamento em $2n$ pontos, e número de caminhos de \texttt{(0, 0)} até \texttt{(n, n)} andando apenas para a direita e para cima, sem passar acima da diagonal.
		\fi

		%% ----------------------------------------------------- %%
	\item \textbf{Número de Narayana.}
		A quantidade de caminhos de $(0, 0)$ a $(N, N)$ que não passam acima da diagonal e com exatamente $M$ picos locais é
		\[
			\frac1N\binom{N}{M}\binom{N}{M - 1}
		\]

		%% ----------------------------------------------------- %%
	\item \textbf{Lema dos ciclos.}
		Dada uma sequência $A_1, A_2, \dots, A_n$ de números tais que $A_i \le 1$ e $K = \sum{A_i} \ge 0$, o número de rotações cíclicas de $A$ tais que cada soma de prefixo é positiva é exatamente $K$.

		%% ----------------------------------------------------- %%
	\item \textbf{Desarranjos.}
		\[
			D_n = n!\sum_{k=0}^{n}\frac{(-1)^k}{k!}
		\]
		e
		\[
			D_n = (n-1)(D_{n-1}+D_{n-2})
		\]

		%% ----------------------------------------------------- %%
	\item \textbf{Funções de parking.}
		O número de sequências $a = (a_1, a_2, \dots, a_m)$ tais que $\forall i, 1 \le a_i \le n$, e para $k = 1, 2, \dots, n$, o número de índices $i$ que satisfazem $a_i \le k$ é pelo menos $m - n + k$ é $(n + 1 - m)(n + 1)^{m - 1}$

		%% ----------------------------------------------------- %%
	\item \textbf{Números de Stirling do segundo tipo.}
		\[
			S(n,k)=S(n-1,k-1)+k\,S(n-1,k)
		\]
		Forma fechada:
		\[
			S(n,k)=\frac1{k!}\sum_{i=0}^{k}(-1)^i\binom{k}{i}(k-i)^n
		\]

		%% ----------------------------------------------------- %%
	\item \textbf{Números de Stirling do primeiro tipo (sem sinal).}
		\[
			c(n,k)=c(n-1,k-1)+(n-1)c(n-1,k)
		\]

		%% ----------------------------------------------------- %%
	\item \textbf{Números de Bell.}
		\[
			B_n=\sum_{k=0}^{n}S(n,k)
		\]

\end{itemize}

%%%% -------------------------------------------------------------------- %%%%
\subsection{Somas e sequências}
\begin{itemize}[leftmargin=*,label={}]
		%% ----------------------------------------------------- %%
		\ifextended
	\item \textbf{Progressão aritmética.}
		\[
			1+2+\cdots+n = \frac{n(n+1)}{2}
		\]
		\[
			a + (a+d) + \cdots + (a+(n-1)d)
			= \frac{n(2a+(n-1)d)}{2}
		\]
		\fi

		%% ----------------------------------------------------- %%
		\ifextended
	\item \textbf{Progressão geométrica.}
		\[
			1+r+r^2+\cdots+r^n = \frac{r^{n+1}-1}{r-1}\qquad (r\ne1)
		\]
		\[
			a+ar+\cdots+ar^{n-1} = a\frac{r^n-1}{r-1}\qquad (r\ne1)
		\]
		\fi

		%% ----------------------------------------------------- %%
	\item \textbf{Somas de potências comuns.}
		\[
			\sum_{i=1}^n i = \frac{n(n+1)}{2}
		\]
		\[
			\sum_{i=1}^n i^2 = \frac{n(n+1)(2n+1)}{6}
		\]
		\[
			\sum_{i=1}^n i^3 = \left(\frac{n(n+1)}{2}\right)^2
		\]

		%% ----------------------------------------------------- %%
	\item \textbf{Convolução binomial (Vandermonde).}
		\[
			\sum_k \binom{a}{k}\binom{b}{n-k}=\binom{a+b}{n}.
		\]

		%% ----------------------------------------------------- %%
		\ifextended
	\item \textbf{Valor esperado por indicadores.}
		Se
		\[
			X=\sum_i I_i,
		\]
		então
		\[
			\mathbb{E}[X]=\sum_i \Pr(I_i=1).
		\]
		\fi

		%% ----------------------------------------------------- %%
		\ifextended
	\item \textbf{Contagem de iterações em submasks.}
		Número total de pares $(m,s)$ com $s\subseteq m\subseteq \{0,\dots,n-1\}$:
		\[
			3^n.
		\]
		\fi

		%% ----------------------------------------------------- %%
		\ifextended
	\item \textbf{Fórmula de contribuição na soma de todos os subarrays.}
		\[
			\sum_{0\le l\le r<n}\ \sum_{i=l}^r a_i
			=
			\sum_{i=0}^{n-1} a_i(i+1)(n-i).
		\]
		O elemento $a_i$ aparece em exatamente $(i+1)(n-i)$ subarrays.
		\fi

		%% ----------------------------------------------------- %%
		\ifextended
	\item \textbf{Contribuição de uma aresta para a soma das distâncias na árvore.}
		Se remover uma aresta de peso $w$ separa a árvore em componentes de tamanhos $s$ e $n-s$, então sua contribuição para
		\[
			\sum_{u<v}\operatorname{dist}(u,v)
		\]
		é
		\[
			w \cdot s(n-s).
		\]
		\fi

		%% ----------------------------------------------------- %%
		\ifextended
	\item \textbf{Matriz de Fibonacci}
		Os números de Fibonacci, definidos por $F_0 = F_1 = 1, F_n = F_{n-1}+F_{n-2}$, podem ser representados pela matriz:
		\[
			\begin{bmatrix}
				1 & 1 \\
				1 & 0
			\end{bmatrix}
		\]
		De tal forma que
		\[
			\begin{bmatrix}
				1 & 1 \\
				1 & 0
			\end{bmatrix}^n = 
			\begin{bmatrix}
				F_n & F_{n-1} \\
				F_{n-1} & F_{n-2}
			\end{bmatrix}
		\]

		%% ----------------------------------------------------- %%
	\item \textbf{Fórmulas para números de Fibonacci}
		\[
			F_{n+m} = F_nF_m + F_{n-1}F_{m-1}
		\]
		\[
			\sum\limits_{i=1}^n F_i = F_{n+2} - 1
		\]
		\[
			\sum\limits_{i=1}^n F_i^2 = F_nF_{n+1}
		\]
		\fi

\end{itemize}

\subsection{Combinatória e Otimização}
\begin{itemize}[leftmargin=*,label={}]
		%% ----------------------------------------------------- %%
	\item \textbf{Teorema de Dilworth.}
		Em qualquer poset finito, o tamanho máximo de uma anticadeia é igual ao número mínimo de cadeias necessárias para cobrir todos os elementos.

		%% ----------------------------------------------------- %%
	\item \textbf{Erdős--Szekeres.}
		Toda sequência de $(r-1)(s-1)+1$ números distintos contém ou uma subsequência crescente de tamanho $r$, ou uma subsequência decrescente de tamanho $s$.

		%% ----------------------------------------------------- %%
	\item \textbf{Teorema de Mirsky.}
		Em qualquer poset finito, o tamanho máximo de uma cadeia é igual ao número mínimo de anticadeias necessárias para cobrir todos os elementos.

		%% ----------------------------------------------------- %%
		\ifextended
	\item \textbf{Mediana minimiza soma de desvios absolutos.}
		\[
			\arg\min_x \sum_i |x-a_i|
		\]
		é qualquer mediana dos $a_i$. Mais geralmente,
		\[
			\arg\min_x \sum_i w_i |x-a_i|
		\]
		é qualquer mediana ponderada.
		\fi

		%% ----------------------------------------------------- %%
		\ifextended
	\item \textbf{Média minimiza soma de quadrados.}
		\[
			\arg\min_x \sum_i (x-a_i)^2
		\]
		é
		\[
			x=\frac{1}{n}\sum_i a_i.
		\]
		Se $x$ precisa ser inteiro, teste $\lfloor \bar a \rfloor$ e $\lceil \bar a \rceil$.
		\fi

		%% ----------------------------------------------------- %%
	\item \textbf{Máximo número de intervalos disjuntos.}
		Ordenar os intervalos por extremo direito crescente e escolher gulosa-mente produz um conjunto máximo de intervalos dois a dois disjuntos.

		%% ----------------------------------------------------- %%
	\item \textbf{Mínimo de pontos para perfurar intervalos.}
		Ordenar os intervalos por extremo direito crescente e escolher repetidamente o extremo direito do primeiro intervalo ainda não coberto produz um conjunto mínimo de pontos que intersecta todos os intervalos.

		%% ----------------------------------------------------- %%
	\item \textbf{Intervalos: mínimo de pontos = máximo de disjuntos.}
		Para uma família de intervalos na reta, o número mínimo de pontos necessários para intersectar todos os intervalos é igual ao número máximo de intervalos dois a dois disjuntos.

		%% ----------------------------------------------------- %%
	\item \textbf{Intervalos dois a dois intersectantes têm interseção comum.}
		Para intervalos na reta, se todo par intersecta, então todos compartilham um ponto em comum. Isso falha em dimensões maiores, mas é muito útil em 1D.

\end{itemize}

\subsection{Teoria dos Números}
\begin{itemize}[leftmargin=*,label={}]
		%% ----------------------------------------------------- %%
	\item \textbf{Truques de floor/ceiling.}
		\[
			\left\lceil \frac{a}{b} \right\rceil = \left\lfloor \frac{a+b-1}{b} \right\rfloor
			\qquad (a,b>0)
		\]
		\[
			\left\lceil \frac{a}{b} \right\rceil = \frac{a+b-1}{b}
		\]
		em aritmética inteira, para $a,b>0$.

		%% ----------------------------------------------------- %%
	\item \textbf{Contagem / soma de divisores.}
		Se
		\[
			n=\prod p_i^{e_i},
		\]
		então
		\[
			d(n)=\prod (e_i+1),
		\]
		\[
			\sigma(n)=\prod \frac{p_i^{e_i+1}-1}{p_i-1}.
		\]

		%% ----------------------------------------------------- %%
		\ifextended
	\item \textbf{Phi de Euler.}
		\[
			\varphi(n)=n\prod_{p\mid n}\left(1-\frac1p\right)
		\]
		\fi

		%% ----------------------------------------------------- %%
	\item \textbf{Triplas pitagóricas.}
		Todas as soluções inteiras primitivas de
		\[
			a^2+b^2=c^2
		\]
		são exatamente
		\[
			a=m^2-n^2,\qquad b=2mn,\qquad c=m^2+n^2,
		\]
		para inteiros $m>n\ge 1$ com
		\[
			\gcd(m,n)=1
			\quad\text{e}\quad
			m\not\equiv n \pmod 2.
		\]
		Todas as soluções inteiras são obtidas multiplicando uma solução primitiva por $k\ge 1$.

		%% ----------------------------------------------------- %%
	\item \textbf{Somas de três quadrados.}
		Um inteiro não negativo $n$ pode ser escrito como
		\[
			n=x^2+y^2+z^2
		\]
		com $x,y,z$ inteiros sse
		\[
			n\ne 4^a(8b+7)
		\]
		para todos os inteiros $a,b\ge 0$.

		%% ----------------------------------------------------- %%
		\ifextended
	\item \textbf{Fórmula de Legendre.}
		Para $p$ primo, o maior expoente de $p$ que divide $n!$ é
		\[
			\sum_{k\ge1}\left\lfloor \frac{n}{p^k}\right\rfloor
		\]
		\fi

		%% ----------------------------------------------------- %%
	\item \textbf{Teorema de Lucas.}
		Para primo $p$, se
		\[
			n = \sum n_i p^i,\qquad k = \sum k_i p^i,
		\]
		então
		\[
			\binom{n}{k} \equiv \prod_i \binom{n_i}{k_i} \pmod p.
		\]

		%% ----------------------------------------------------- %%
	\item \textbf{Teorema de Kummer.}
		O expoente de um primo $p$ em $\binom{n}{k}$ é igual ao número de vai-uns ao somar $k$ e $n-k$ na base $p$.

		%% ----------------------------------------------------- %%
	\item \textbf{LTE (versão comum).}
		Se $p$ é um primo ímpar, $p \mid (a-b)$ e $p \nmid ab$, então
		\[
			v_p(a^n-b^n)=v_p(a-b)+v_p(n).
		\]
		Além disso, para $n$ ímpar e $p \mid (a+b)$,
		\[
			v_p(a^n+b^n)=v_p(a+b)+v_p(n).
		\]

		%% ----------------------------------------------------- %%
		\ifextended
	\item \textbf{Condição geral de compatibilidade do CRT.}
		O sistema
		\[
			x \equiv a_i \pmod{m_i}
		\]
		tem solução sse, para todos $i,j$,
		\[
			a_i \equiv a_j \pmod{\gcd(m_i,m_j)}.
		\]
		Quando tem solução, ela é única módulo $\mathrm{lcm}(m_1,\dots,m_k)$.
		\fi

		%% ----------------------------------------------------- %%
	\item \textbf{Inversão de Möbius.}
		Se
		\[
			g(n)=\sum_{d \mid n} f(d),
		\]
		então
		\[
			f(n)=\sum_{d \mid n} \mu(d)\, g\!\left(\frac{n}{d}\right).
		\]

		%% ----------------------------------------------------- %%
		\ifextended
	\item \textbf{Truque de agrupar quocientes.}
		O valor de
		\[
			\left\lfloor \frac{n}{i}\right\rfloor
		\]
		muda apenas $O(\sqrt n)$ vezes quando $i$ varia. Mais precisamente, todos os $i$ com o mesmo quociente formam um intervalo
		\[
			i \in \left[\left\lfloor \frac{n}{q+1}\right\rfloor+1,\ \left\lfloor \frac{n}{q}\right\rfloor\right].
		\]
		\fi

		%% ----------------------------------------------------- %%
	\item \textbf{Existência de raiz primitiva.}
		Existem raízes primitivas módulo $n$ sse
		\[
			n \in \{1,2,4,p^k,2p^k\},
		\]
		onde $p$ é um primo ímpar.

		%% ----------------------------------------------------- %%
		\ifextended
	\item \textbf{Teorema de Wilson.}
		Para primo $p$,
		\[
			(p-1)! \equiv -1 \pmod p.
		\]
		\fi
\end{itemize}

\subsection{Grafos e Jogos}
\begin{itemize}[leftmargin=*,label={}]
		%% ----------------------------------------------------- %%
	\item \textbf{Teorema do Casamento de Hall.}
		Um grafo bipartido $(L,R)$ tem um matching saturando $L$ sse
		\[
			\forall S \subseteq L,\quad |N(S)| \ge |S|.
		\]

		%% ----------------------------------------------------- %%
	\item \textbf{Teorema de König.}
		Em um grafo bipartido,
		\[
			\text{matching máximo} = \text{cobertura mínima}.
		\]

		%% ----------------------------------------------------- %%
	\item \textbf{Condições para trilha/circuito euleriano (direcionado).}
		Ignorando os vértices isolados, todos os vértices relevantes devem pertencer a uma única componente fracamente conexa. Então:
		\begin{itemize}
			\item existe circuito euleriano sse $\deg^+(v)=\deg^-(v)$ para todo vértice;
			\item existe trilha euleriana sse exatamente um vértice tem $\deg^+(v)=\deg^-(v)+1$, exatamente um tem $\deg^-(v)=\deg^+(v)+1$, e todos os outros são balanceados.
		\end{itemize}

		%% ----------------------------------------------------- %%
	\item \textbf{Condições para trilha/circuito euleriano (não direcionado).}
		Ignorando os vértices isolados, todos os vértices relevantes devem pertencer a uma única componente conexa. Então:
		\begin{itemize}
			\item existe circuito euleriano sse todo vértice tem grau par;
			\item existe trilha euleriana sse exatamente $0$ ou $2$ vértices têm grau ímpar.
		\end{itemize}

		%% ----------------------------------------------------- %%
	\item \textbf{Critério de 2-SAT.}
		Uma instância de 2-SAT é satisfatível sse, para toda variável $x$, os literais $x$ e $\neg x$ estão em SCCs diferentes do grafo de implicação.

		%% ----------------------------------------------------- %%
	\item \textbf{Restrições de diferença $\leftrightarrow$ menores caminhos.}
		Um sistema de desigualdades da forma
		\[
			x_v - x_u \le w
		\]
		pode ser modelado por uma aresta $u \to v$ de peso $w$. Então:
		\begin{itemize}
			\item o sistema é viável sse não existe ciclo negativo;
			\item uma solução viável é dada pelas distâncias a partir de uma super-fonte.
		\end{itemize}

		%% ----------------------------------------------------- %%
	\item \textbf{Cobertura por caminhos em DAG via matching.}
		O número mínimo de caminhos disjuntos em vértices que cobrem todos os vértices de um DAG é
		\[
			n - \text{(matching máximo na expansão bipartida do DAG)}.
		\]

		%% ----------------------------------------------------- %%
	\item \textbf{Teorema Matrix-Tree de Kirchhoff.}
		O número de árvores geradoras de um grafo é igual a qualquer cofator de sua matriz Laplaciana. (A matriz de graus - a matriz de adjacência).

		%% ----------------------------------------------------- %%
		\ifextended
	\item \textbf{Teorema de Sprague--Grundy.}
		Todo jogo imparcial finito em normal play é equivalente a uma pilha de Nim cujo tamanho é seu número de Grundy; o número de Grundy é o mex dos números de Grundy das jogadas possíveis. O número de Grundy de uma soma disjunta é o xor dos números de Grundy das componentes.
		\fi

		%% ----------------------------------------------------- %%
	\item \textbf{Jogo UGV.}
		Considere o jogo de mover ao longo de arestas de um grafo sem poder repetir vértices, e o jogador que não puder mover perde. O segundo jogador tem uma estratégia vencedora sse existe um matching máximo no grafo que deixa o vértice inicial insaturado.

\end{itemize}

\subsection{Strings}
\begin{itemize}[leftmargin=*,label={}]
		%% ----------------------------------------------------- %%
	\item \textbf{Teorema de Fine--Wilf.}
		Se uma string tem períodos $p$ e $q$, e seu tamanho é pelo menos
		\[
			p+q-\gcd(p,q),
		\]
		então ela também tem período $\gcd(p,q)$.

		%% ----------------------------------------------------- %%
		\ifextended
	\item \textbf{Hash polinomial de substring.}
		Se
		\[
			H[i+1]=H[i]\cdot B+s[i],
		\]
		então
		\[
			hash(l,r)=H[r]-H[l]\cdot B^{r-l}
		\]
		para o intervalo semiaberto $[l,r)$.
		\fi

		%% ----------------------------------------------------- %%
	\item \textbf{Substrings distintas.}
		Usando array de sufixos:
		\[
			\frac{n(n+1)}{2} - \sum lcp[i]
		\]
		Usando autômato de sufixos:
		\[
			\sum len(i) - len(link(i))
		\]

		%% ----------------------------------------------------- %%
	\item \textbf{Notas de automato de sufixos.}
		O automato de sufixos é um DAG. Cada caminho a partir da raiz representa uma substring. Cada nó representa um conjunto de \texttt{endpos}. Para obter o tamanho dos conjuntos de endpos, basta propagar os \texttt{cnt} por suffix links na ordem contrária de \texttt{ord}. \texttt{link(i)} é o maior sufixo com um conjunto \texttt{endpos} diferente.

		%% ----------------------------------------------------- %%
	\item \textbf{Encontrando a maior substring comum.}
		Construa o automato de sufixos de uma das strings. Itere pela segunda string mantendo o maior sufixo dela que aparece na primeira string. Quando uma transição falha, basta subir por suffix links, e o novo maior tamanho do sufixo vira \texttt{len(v)}.

\end{itemize}

\subsection{Geometria}
\begin{itemize}[leftmargin=*,label={}]
		%% ----------------------------------------------------- %%
	\item \textbf{Transformações da distância Manhattan.}
		Em 2D,
		\ifextended
		\[
			|x_1-x_2|+|y_1-y_2|
			= \max\Bigl(
			|(x+y)'-(x+y)|,\,
			|(x-y)'-(x-y)|
			\Bigr)
		\]
		Na prática, para a maior distância entre pares:
		\[
			\max_{i,j}|x_i-x_j|+|y_i-y_j| = \max(D_+, D_-)
		\]
		onde
		\[
			D_+ = \max(x+y)-\min(x+y),\qquad
			D_- = \max(x-y)-\min(x-y)
		\]
		\else
		\[
			L_1 = \max(|\Delta_+|, |\Delta_-|)
		\]
		\[
			\Delta_\pm = (x\pm y)'-(x\pm y)
		\]
		Na prática:
		\[
			\max L_1 = \max(D_+, D_-)
		\]
		\[
			D_\pm = \max(x\pm y)-\min(x\pm y)
		\]
		\fi

		%% ----------------------------------------------------- %%
	\item \textbf{Teorema de Euler para grafos planares.}
		Num grafo planar com $v$ vértices, $e$ arestas, $f$ faces e $k$ componentes conexas,
		\[
			v - e + f = 1 + k
		\]
		%% ----------------------------------------------------- %%
	\item \textbf{Teorema de Pick.}
		Para um polígono simples na grid,
		\[
			A = I + \frac{B}{2} - 1,
		\]
		onde $I$ é o número de pontos de grid internos e $B$ é o número de pontos de grid na fronteira.

		%% ----------------------------------------------------- %%
	\item \textbf{Pontos de grid em um segmento.}
		O número de pontos de grid no segmento de $(x_1,y_1)$ até $(x_2,y_2)$ é
		\[
			\gcd(|x_1-x_2|, |y_1-y_2|) + 1.
		\]

		%% ----------------------------------------------------- %%
		\ifextended
	\item \textbf{Distância de ponto a reta.}
		Distância do ponto $P$ à reta passando por $A,B$:
		\[
			\frac{|\operatorname{cross}(B-A,P-A)|}{|B-A|}
		\]
		\fi

		%% ----------------------------------------------------- %%
	\item \textbf{Fórmula de Heron.}
		Para lados $a,b,c$ e semiperímetro
		\[
			s=\frac{a+b+c}{2},
		\]
		a área do triângulo é
		\[
			\sqrt{s(s-a)(s-b)(s-c)}.
		\]

		%% ----------------------------------------------------- %%
		\ifextended
	\item \textbf{Lei dos cossenos.}
		\[
			c^2=a^2+b^2-2ab\cos C
		\]
		\fi

		%% ----------------------------------------------------- %%
		\ifextended
	\item \textbf{Ângulo pelo produto escalar.}
		\[
			\cos\theta = \frac{a\cdot b}{|a||b|}
		\]
		\fi

		%% ----------------------------------------------------- %%
	\item \textbf{Raio da circunferência circunscrita e inscrita.}
		Para área do triângulo $\Delta$:
		\[
			R=\frac{abc}{4\Delta},\qquad r=\frac{\Delta}{s}
		\]
		onde $s=\frac{a+b+c}{2}$.

		%% ----------------------------------------------------- %%
		\ifextended
	\item \textbf{Comprimento da corda.}
		Para raio $r$ e ângulo central $\theta$:
		\[
			2r\sin\frac{\theta}{2}
		\]
		\fi

		%% ----------------------------------------------------- %%
		\ifextended
	\item \textbf{Teorema de Helly em $\mathbb{R}^2$.}
		Para uma família finita de conjuntos convexos no plano, se todo subconjunto de $3$ deles intersecta, então todos intersectam.
		\fi

		%% ----------------------------------------------------- %%
	\item \textbf{Teorema de Helly em $\mathbb{R}^d$.}
		Para conjuntos convexos em $\mathbb{R}^d$, se todo subconjunto de $d+1$ deles intersecta, então todos intersectam.

		%% ----------------------------------------------------- %%
		\ifextended
	\item \textbf{Teorema de Carathéodory.}
		Se um ponto pertence ao fecho convexo de um conjunto em $\mathbb{R}^d$, então ele pertence ao fecho convexo de no máximo $d+1$ pontos desse conjunto.
		\fi
\end{itemize}

\subsection{DP}
\begin{itemize}[leftmargin=*,label={}]
		%% ----------------------------------------------------- %%
	\item \textbf{Formular DP}
		Se um problema pode ser resolvido com DP, mesmo que a solução óbvia resultaria em TLE, escreva a formulação da DP. Talvez seja possível otimizar a DP e resolver rápido o suficiente, ou mesmo generalizar a forma obtida para algo mais simples/fácil de calcular.

		%% ----------------------------------------------------- %%
	\item \textbf{Otimizar com matrizes}
		Especialmente útil quando o problema pode ser resolvido com uma DP que depende apenas de uma janela de estados anteriores de tamanho fixo. Tentar representar as transições da DP como multiplicação de matrizes.

		%% ----------------------------------------------------- %%
	\item \textbf{Otimizar com funções}
		Pode ser que a DP tome o formato de uma função. Neste caso, talvez seja possível calcular o valor da DP em outros pontos descobrindo a forma da função, como usando interpolação ou outra forma.

		%% ----------------------------------------------------- %%
		\ifextended
	\item \textbf{Desigualdade de Monge.}
		Um array $A$ é Monge se
		\[
			A[i][j]+A[i'][j'] \le A[i][j']+A[i'][j]
		\]
		para $i\le i'$ e $j\le j'$. Essa é a condição estrutural por trás de várias otimizações de DP.
		\fi

		%% ----------------------------------------------------- %%
		\ifextended
	\item \textbf{Condição da otimização de Knuth.}
		Para DP de intervalos,
		\[
			opt[l][r-1] \le opt[l][r] \le opt[l+1][r]
		\]
		sob as hipóteses usuais de desigualdade do quadrilátero / monotonicidade.
		\fi

		%% ----------------------------------------------------- %%
	\item \textbf{DP contando subsequências}
		\texttt{DP[i]} = número de subsequências distintas até a posição \texttt{i}.
		Temos que:
		\[
			DP[i] = 2 * DP[i - 1] - DP[j - 1]
		\]
		onde \texttt{j} é a última posição que o valor na posição \texttt{i} apareceu.

		%% ----------------------------------------------------- %%
	\item \textbf{Convolução Min-Plus.}
		A convolução Min-Plus de duas sequências convexas pode ser calculada com greedy nas sequências de diferença.
		A convolução Min-Plus de duas sequências convexas também é convexa.

		%% ----------------------------------------------------- %%
		\ifextended
	\item \textbf{Truque dos intervalos circulares.}
		Em problemas sobre um círculo, cortar em um ponto escolhido transforma intervalos circulares em intervalos comuns. Muitas vezes isso se resolve duplicando o array/lista de ângulos e trabalhando numa estrutura linearizada de tamanho $2n$.
		\fi

\end{itemize}

\subsection{Tabela de funções geradoras}
\begingroup
\allowdisplaybreaks[4]
\begin{alignat*}{2}
	[x^n]F(x) &\qquad && F(x) \\
	\rlap{\hspace{-0.18\linewidth}\rule{0.45\linewidth}{0.4pt}} &&& \\[-1.4ex]
	1 &\qquad && \dfrac{1}{1 - x} \\[6pt]
	a^n &\qquad && \dfrac{1}{1 - ax} \\[6pt]
	n &\qquad && \dfrac{x}{(1 - x)^2} \\[6pt]
	\binom{n + k}{k} &\qquad && \dfrac{1}{(1 - x)^{k + 1}} \\[6pt]
	\binom{n - a + k - 1}{k - 1} &\qquad && \dfrac{x^a}{(1-x)^{k}} \\[6pt]
	\binom{k}{n} &\qquad && (1 + x)^k \\[6pt]
	C_n &\qquad && \dfrac{1 - \sqrt{1 - 4x}}{2x} \\[6pt]
	F_n &\qquad && \dfrac{x}{1 - x - x^2} \\[6pt]
	\dfrac{D_n}{n!} &\qquad && \dfrac{e^{-x}}{1 - x} \\[6pt]
	\dfrac{1}{n} &\qquad && \log\dfrac{1}{1-x}
\end{alignat*}
\endgroup
